%%%%%%%%%%%%%%%%%%%%%%%%%%%%%%%%%%%%%%%%%
% Boletim semanal estadual do InfoDengue
% Adapted for InfoDengue by Claudia Codeco and Thais Riback (May 2016)
%
% Created by:
% Bob Kerstetter (https://www.tug.org/texshowcase/) and extensively modified by:
% Vel (vel@latextemplates.com)
% 
% Original template downloaded from:
% http://www.LaTeXTemplates.com
%
% License:
% CC BY-NC-SA 3.0 (http://creativecommons.org/licenses/by-nc-sa/3.0/)
%%%%%%%%%%%%%%%%%%%%%%%%%%%%%%%%%%%%%%%%%

\documentclass[10pt]{article} % The default font size is 10pt; 11pt and 12pt are alternatives
%%%%%%%%%%%%%%%%%%%%%%%%%%%%%%%%%%%%%%%%%
% Professional Newsletter Template
% Structural Definitions File
% Version 1.0 (09/03/14)
%
% Created by:
% Vel (vel@latextemplates.com)
% 
% This file has been downloaded from:
% http://www.LaTeXTemplates.com
%
% License:
% CC BY-NC-SA 3.0 (http://creativecommons.org/licenses/by-nc-sa/3.0/)
%
%%%%%%%%%%%%%%%%%%%%%%%%%%%%%%%%%%%%%%%%%

%----------------------------------------------------------------------------------------
%	REQUIRED PACKAGES
%----------------------------------------------------------------------------------------

\usepackage{graphicx} % Required for including images
\usepackage[utf8]{inputenc}
\usepackage{microtype} % Improved typography
\usepackage{multicol} % Used for the two-column layout of the document
\usepackage{booktabs} % Required for nice horizontal rules in tables
\usepackage{wrapfig} % Required for in-line images
\usepackage{float} % Required for forcing figures not to float with the [H] 
\usepackage{longtable}
\usepackage{fancyhdr} % footer 
%parameter
\usepackage{Sweave}

%-------------

% -------------
% footer style
\pagestyle{fancy}
\fancyhead{}
\fancyfoot{}
\renewcommand{\headrulewidth}{0.0pt}
\renewcommand{\footrulewidth}{0.4pt}

%------------------------------------------------
% Fonts

\usepackage{charter} % Use the Charter font as the main document font
\usepackage{courier} % Use the Courier font for \texttt (monospaced) only
%\usepackage[T1]{fontenc} % Use T1 font encoding
\usepackage[utf8]{inputenc}
\usepackage[brazilian]{babel}

%------------------------------------------------
% List Separation

\usepackage{enumitem} % Required to customize the list environments
\setlist{noitemsep,nolistsep} % Remove spacing before, after and within lists for a compact look

%------------------------------------------------
% Figure and Table Caption Styles

\usepackage{caption} % Required for changing caption styles
\captionsetup[table]{labelfont={bf,sf},labelsep=period,justification=justified} % Specify the table caption style
\captionsetup[figure]{labelfont={sf,bf},labelsep=period,justification=justified, font=small} % Specify the figure caption style
\setlength{\abovecaptionskip}{10pt} % Whitespace above captions

%------------------------------------------------
% Spacing Between Paragraphs

\makeatletter
\usepackage{parskip}
\setlength{\parskip}{6pt}
\newcommand{\@minipagerestore}{\setlength{\parskip}{6pt}}
\makeatother

%----------------------------------------------------------------------------------------
%	PAGE MARGINS AND SPACINGS
%----------------------------------------------------------------------------------------

\textwidth = 7 in % Text width
\textheight = 9.5 in % Text height
\oddsidemargin = -10pt % Left side margin on odd pages
\evensidemargin = -10pt % Left side margin on even pages
\topmargin = -10pt % Top margin
\headheight = 0pt % Remove the header by setting its space to 0
\headsep = 0pt % Remove the space between the header and top of the page
\parskip = 4pt % Space between paragraph
\parindent = 0.0in % Paragraph indentation
%\pagestyle{empty} % Disable page numbering

%----------------------------------------------------------------------------------------
%	COLORS
%----------------------------------------------------------------------------------------

\usepackage[dvipsnames,svgnames]{xcolor} % Required to specify custom colors
\definecolor{altncolor}{rgb}{.0,0.0,0.8} % Dark blue
\definecolor{myred}{rgb}{.8,0,0} % Dark red
\definecolor{myorange}{rgb}{.93,.57,.0} % Dark orange
\definecolor{myyellow}{rgb}{.93,.93,.0} % yellow
\definecolor{mygreen}{rgb}{.0,.55,.27} % dark green
\definecolor{mywhite}{rgb}{1,1,1} % dark green

\usepackage[colorlinks=true, linkcolor=altncolor, anchorcolor=altncolor, citecolor=altncolor, filecolor=altncolor, menucolor=altncolor, urlcolor=altncolor]{hyperref} % Use the color defined above for all links

% --------------------------------------
% Bullet styles
\usepackage{pifont}
\newcommand{\gsquare}{\item[\color{mygreen}\ding{110}]} 
\newcommand{\ysquare}{\item[\color{myyellow}\ding{110}]} 
\newcommand{\osquare}{\item[\color{myorange}\ding{110}]} 
\newcommand{\rsquare}{\item[\color{myred}\ding{110}]} 
\newcommand{\wsquare}{\item[\color{mywhite}\ding{110}]} 


%----------------------------------------------------------------------------------------
%	BOX STYLES
%----------------------------------------------------------------------------------------

\usepackage[framemethod=TikZ]{mdframed}% Required for creating boxes
\mdfdefinestyle{sidebar}{
    linecolor=black, % Outer line color
    outerlinewidth=0.5pt, % Outer line width
    roundcorner=0pt, % Amount of corner rounding
    innertopmargin=10pt, % Top margin
    innerbottommargin=10pt, % Bottom margin
    innerrightmargin=10pt, % Right margin
    innerleftmargin=10pt, % Left margin
    backgroundcolor=white, % Box background color
    frametitlebackgroundcolor=white, % Title background color
    frametitlerule=false, % Title rule - true or false
    frametitlerulecolor=white, % Title rule color
    frametitlerulewidth=0.5pt, % Title rule width
    frametitlefont=\Large, % Title heading font specification
    font=\small
}

\mdfdefinestyle{intextbox}{
    linecolor=blue, % Outer line color
    outerlinewidth=0.5pt, % Outer line width
    roundcorner=10pt, % Amount of corner rounding
    innertopmargin=7pt, % Top margin
    innerbottommargin=7pt, % Bottom margin
    innerrightmargin=7pt, % Right margin
    innerleftmargin=7pt, % Left margin
    backgroundcolor=white, % Box background color
    frametitlebackgroundcolor=white, % Title background color
    frametitlerule=false, % Title rule - true or false
    frametitlerulecolor=white, % Title rule color
    frametitlerulewidth=0.5pt, % Title rule width
    frametitlefont=\Large % Title heading font specification
}

%----------------------------------------------------------------------------------------
%	HEADING STYLE
%----------------------------------------------------------------------------------------

\newcommand{\heading}[2]{ % Define the \heading command
%\vspace{#2} % White space above the heading
{\begin{center}\Large\textbf{#1}\end{center}} % The heading style
%\vspace{#2} % White space below the heading
}

%\newcommand{\Inicio}{\hyperlink{contents}{{\small Inicio}}} % Define a command for linking back to the contents of the newsletter
\newcommand{\BackToContents}{\hyperlink{contents}{{\small Início}}} % Include the document which specifies all packages and structural customizations for this template

% Carregando dados  (atualmente é manual para cada estado)

\fancyfoot[C]{Boletim Estadual - Rio de Janeiro}
\fancyfoot[R]{Semana 43 de 2016}
\fancyfoot[L]{\href{http://info.dengue.mat.br}{InfoDengue}}

\usepackage{Sweave}
\begin{document}
\Sconcordance{concordance:BoletimEstadual_RJ.tex:BoletimEstadual_RJ.Rnw:%
1 18 1 1 7 5 1 1 0 33 1 1 6 7 0 1 4 67 1 1 64 172 0 1 2 5 1 1 4 66 1}


%---------------------------------------------------------------------------------
%	HEADER IMAGE
%---------------------------------------------------------------------------------

\begin{figure}[H]
\centering\includegraphics[width=1\linewidth]{../reportconfig/InfoDengue2.png}  
\end{figure}

\centerline {\color{altncolor}\rule{1\linewidth}{2.75pt}} % Horizontal line

%---------------------------------------------------------------------------------
%	CAIXA LATERAL - PRIMEIRA PAGINA
%--------------------------------------------------------------------------------

\begin{minipage}[t]{.30\linewidth} % Mini page taking up 30% of the actual page
\begin{mdframed}[style=sidebar,frametitle={}] % Sidebar box

%-----------------------------------------------------------

\hypertarget{contents}{\textbf{{\large Boletim Semanal}}} 

\textbf{Semana 43 de 2016} % se, ano, em pp.RData

\begin{itemize}
\item \hyperlink{estado}{O Estado} 
\end{itemize}


\textbf{Regionais de Saúde:}


\begin{itemize}\item \hyperlink{LtrS}{Litoral Sul}
\item \hyperlink{Nrst}{Noroeste}
\item \hyperlink{BxdL}{Baixada Litorânea}
\item \hyperlink{MdPr}{Médio Paraíba}
\item \hyperlink{MtrI}{Metropolitana I}
\item \hyperlink{Srrn}{Serrana}
\item \hyperlink{Nort}{Norte}
\item \hyperlink{CntS}{Centro Sul}
\item \hyperlink{MtII}{Metropolitana II}\end{itemize}
\textbf{Municípios:}

\begin{itemize}
\item \hyperlink{municips}{Tabela de Municípios} 
\end{itemize}

\centerline {\rule{.75\linewidth}{.25pt}} % Horizontal line

%-----------------------------------------------------------

\hyperlink{vartab}{Variáveis nas Tabelas}

\hyperlink{notas}{Notas} % These link to their appropriate sections in the newsletter

\hyperlink{creditos}{Creditos} % These link to their appropriate sections in the newsletter

\centerline {\rule{.75\linewidth}{.25pt}} % Horizontal line

\textbf{Contato}
\begin{description}
\item \href{mailto:alerta\_dengue@fiocruz.br}{alerta\_dengue@fiocruz.br}  
\end{description}

\end{mdframed}
\end{minipage}\hfill % End the sidebar mini page 
\begin{minipage}[t]{.66\linewidth} % Nao pode colocar espaco acima senao ela nao fica lado a lado

%--------------------------------------------------------------------------------
%	TEXTO PRINCIPAL - PRIMEIRA PAGINA - ALERTA A NIVEL DO ESTADO
%-------------------------------------------------------------------------------
\hypertarget{estado}{\heading{Situação da Dengue no Estado do Rio de Janeiro}{6pt}} % \hypertarget provides a label to reference using 

Desde o início do ano, 108014 casos foram registrados no estado, sendo 46 na última semana. A figura abaixo mostra as condições de transmissão em cada município.

% Mapa estadual
\includegraphics[width=0.8\textwidth]{figs/Mapa_ERJ.png}
%-----------------------------------------------------------

Dos 92 municipios, 73 encontram-se em nivel verde, 19 em nivel amarelo, 0 em nivel laranja e 0 em nivel vermelho referentes a semana epidemiológica 43-2016. Para informações mais atualizadas
sobre o município do Rio de Janeiro ter acesso ao mapa interativo do estado consultar em \href{http://info.dengue.mat.br}{\textit{Info Dengue}}.


%--------------------------------------------------------------------------------
%	PRIMEIRA PAGINA : BOX - CODIGO DE CORES
%--------------------------------------------------------------------------------
\vspace{1cm}
\begin{mdframed}[style=intextbox,frametitle={}] % Sidebar box

\hypertarget{descriptivebox}{\heading{O código de Cores}{1pt}} % \hypertarget provides a label to reference using \hyperlink{label}{link text}
As cores indicam niveis de atenção
\begin{description}
\item[Verde:] temperaturas amenas, baixa incidência de casos.      
\item[Amarelo:] temperatura propícia para a população do vetor e transmissão da dengue.
\item[Laranja:] transmissão aumentada e sustentada de dengue. 
\item[Vermelho:] incidência alta de dengue, acima dos 90\% históricos.
\end{description}
\end{mdframed}

 \textbf{Na semana passada:} 19 municípios em nivel amarelo, 0 em laranja e 0 em vermelho.    

\end{minipage} % End the main body - first page mini page


%----------------------------------------------------------------------------------
%	SEGUNDA PAGINA: TEXTO PRINCIPAL - LOOP PELAS REGIONAIS
%----------------------------------------------------------------------------------

\begin{minipage}[t]{.66\linewidth}
\hypertarget{LtrS}{\heading{Regional Litoral Sul}{6pt}}
\includegraphics[width=0.8\textwidth]{figs/MapaRJ_LitoralSul.png}\vspace{0.5cm}\captionof{table}{Resumo das últimas seis semanas epidemiológicas na Regional Litoral Sul }\begin{center}
% latex table generated in R 3.3.0 by xtable 1.8-2 package
% Tue Aug 30 11:26:27 2016
\begin{tabular}{c|ccccccc}
  \hline
SE & temperatura & tweet & casos notif & casos preditos & ICmin & ICmax & incidência \\ 
  \hline
201629 & 12 & 0 & 4 & 4 & 4 & 4 & 2 \\ 
  201630 & 14 & 0 & 9 & 9 & 9 & 9 & 3 \\ 
  201631 & 16 & 0 & 6 & 6 & 6 & 6 & 2 \\ 
  201632 & 14 & 2 & 0 & 0 & 0 & 0 & 0 \\ 
  201633 & 16 & 0 & 0 & 0 & 0 & 0 & 0 \\ 
  201634 & 14 & 0 & 0 & 0 & 0 & 0 & 0 \\ 
   \hline
\end{tabular}

\end{center}
\small{\hyperlink{vartab}{ver descrição das variáveis}}\vspace{0.5cm}\begin{center}
\captionof{figure}{Casos notificados de dengue e Índice de menção em midia social sobre dengue na Regional Litoral Sul }\includegraphics[width=1.3\textwidth]{figs/tweetRJ_LitoralSul.png}\end{center}
\end{minipage}\hfill\begin{minipage}[t]{.30\linewidth}
\begin{mdframed}[style=sidebar,frametitle={}]
\textbf{\hyperlink{municips}{Municipios}}\begin{itemize}\gsquare Angra dos Reis 
\ysquare Mangaratiba 
\gsquare Paraty 
\end{itemize}\BackToContents\end{mdframed}\hfill\end{minipage}\newpage\begin{minipage}[t]{.66\linewidth}
\hypertarget{Nrst}{\heading{Regional Noroeste}{6pt}}
\includegraphics[width=0.8\textwidth]{figs/MapaRJ_Noroeste.png}\vspace{0.5cm}\captionof{table}{Resumo das últimas seis semanas epidemiológicas na Regional Noroeste }\begin{center}
% latex table generated in R 3.3.1 by xtable 1.8-2 package
% Tue Oct 18 10:05:44 2016
\begin{tabular}{c|ccccccc}
  \hline
SE & temperatura & tweet & casos notif & casos preditos & ICmin & ICmax & incidência \\ 
  \hline
201636 & 22 & 1 & 1 & 1 & 1 & 1 & 0 \\ 
  201637 & 22 & 0 & 2 & 2 & 2 & 2 & 1 \\ 
  201638 & 21 & 0 & 0 & 0 & 0 & 0 & 0 \\ 
  201639 & 21 & 0 & 3 & 3 & 3 & 3 & 1 \\ 
  201640 &  & 0 & 4 & 4 & 4 & 5 & 1 \\ 
  201641 &  &  & 0 & 0 & 0 & 0 & 0 \\ 
   \hline
\end{tabular}

\end{center}
\small{\hyperlink{vartab}{ver descrição das variáveis}}\vspace{0.5cm}\begin{center}
\captionof{figure}{Casos notificados de dengue e Índice de menção em midia social sobre dengue na Regional Noroeste }\includegraphics[width=1.3\textwidth]{figs/tweetRJ_Noroeste.png}\end{center}
\end{minipage}\hfill\begin{minipage}[t]{.30\linewidth}
\begin{mdframed}[style=sidebar,frametitle={}]
\textbf{\hyperlink{municips}{Municipios}}\begin{itemize}\ysquare Aperibé 
\ysquare Italva 
\ysquare Itaperuna 
\ysquare Laje do Muriaé 
\ysquare Miracema 
\ysquare Natividade 
\ysquare Porciúncula 
\ysquare Varre-Sai 
\ysquare Santo Antônio de Pádua 
\ysquare São José de Ubá 
\ysquare Bom Jesus do Itabapoana 
\ysquare Cambuci 
\ysquare Cardoso Moreira 
\gsquare Itaocara 
\end{itemize}\BackToContents\end{mdframed}\hfill\end{minipage}\newpage\begin{minipage}[t]{.66\linewidth}
\hypertarget{BxdL}{\heading{Regional Baixada Litorânea}{6pt}}
\includegraphics[width=0.8\textwidth]{figs/MapaRJ_BaixadaLitoranea.png}\vspace{0.5cm}\captionof{table}{Resumo das últimas seis semanas epidemiológicas na Regional Baixada Litorânea }\begin{center}
% latex table generated in R 3.3.1 by xtable 1.8-2 package
% Tue Sep 20 20:49:28 2016
\begin{tabular}{c|ccccccc}
  \hline
SE & temperatura & tweet & casos notif & casos preditos & ICmin & ICmax & incidência \\ 
  \hline
201632 & 19 & 5 & 20 & 20 & 20 & 20 & 3 \\ 
  201633 & 20 & 4 & 9 & 9 & 9 & 9 & 1 \\ 
  201634 & 17 & 1 & 13 & 13 & 13 & 13 & 2 \\ 
  201635 & 20 &  & 7 & 7 & 7 & 7 & 1 \\ 
  201636 & 21 &  & 7 & 7 & 7 & 7 & 1 \\ 
  201637 & 21 &  & 9 & 11 & 9 & 12 & 1 \\ 
   \hline
\end{tabular}

\end{center}
\small{\hyperlink{vartab}{ver descrição das variáveis}}\vspace{0.5cm}\begin{center}
\captionof{figure}{Casos notificados de dengue e Índice de menção em midia social sobre dengue na Regional Baixada Litorânea }\includegraphics[width=1.3\textwidth]{figs/tweetRJ_BaixadaLitoranea.png}\end{center}
\end{minipage}\hfill\begin{minipage}[t]{.30\linewidth}
\begin{mdframed}[style=sidebar,frametitle={}]
\textbf{\hyperlink{municips}{Municipios}}\begin{itemize}\gsquare Araruama 
\gsquare Armação dos Búzios 
\gsquare Arraial do Cabo 
\gsquare Cabo Frio 
\gsquare Casimiro de Abreu 
\gsquare São Pedro da Aldeia 
\gsquare Saquarema 
\gsquare Rio das Ostras 
\gsquare Iguaba Grande 
\end{itemize}\BackToContents\end{mdframed}\hfill\end{minipage}\newpage\begin{minipage}[t]{.66\linewidth}
\hypertarget{MdPr}{\heading{Regional Médio Paraíba}{6pt}}
\includegraphics[width=0.8\textwidth]{figs/MapaRJ_MedioParaiba.png}\vspace{0.5cm}\captionof{table}{Resumo das últimas seis semanas epidemiológicas na Regional Médio Paraíba }\begin{center}
% latex table generated in R 3.3.1 by xtable 1.8-2 package
% Tue Oct  4 11:51:46 2016
\begin{tabular}{c|ccccccc}
  \hline
SE & temperatura & tweet & casos notif & casos preditos & ICmin & ICmax & incidência \\ 
  \hline
201634 & 15 & 1 & 85 & 88 & 85 & 89 & 10 \\ 
  201635 & 18 & 0 & 108 & 115 & 112 & 116 & 12 \\ 
  201636 & 19 & 2 & 5 & 5 & 5 & 5 & 1 \\ 
  201637 & 20 & 1 & 95 & 114 & 108 & 116 & 11 \\ 
  201638 &  & 0 & 6 & 6 & 6 & 7 & 1 \\ 
  201639 &  & 1 & 32 & 62 & 55 & 65 & 4 \\ 
   \hline
\end{tabular}

\end{center}
\small{\hyperlink{vartab}{ver descrição das variáveis}}\vspace{0.5cm}\begin{center}
\captionof{figure}{Casos notificados de dengue e Índice de menção em midia social sobre dengue na Regional Médio Paraíba }\includegraphics[width=1.3\textwidth]{figs/tweetRJ_MedioParaiba.png}\end{center}
\end{minipage}\hfill\begin{minipage}[t]{.30\linewidth}
\begin{mdframed}[style=sidebar,frametitle={}]
\textbf{\hyperlink{municips}{Municipios}}\begin{itemize}\gsquare Barra Mansa 
\gsquare Itatiaia 
\gsquare Pinheiral 
\gsquare Porto Real 
\gsquare Rio Claro 
\gsquare Rio das Flores 
\gsquare Piraí 
\gsquare Quatis 
\gsquare Resende 
\gsquare Valença 
\gsquare Volta Redonda 
\gsquare Barra do Piraí 
\end{itemize}\BackToContents\end{mdframed}\hfill\end{minipage}\newpage\begin{minipage}[t]{.66\linewidth}
\hypertarget{MtrI}{\heading{Regional Metropolitana I}{6pt}}
\includegraphics[width=0.8\textwidth]{figs/MapaRJ_MetropolitanaI.png}\vspace{0.5cm}\captionof{table}{Resumo das últimas seis semanas epidemiológicas na Regional Metropolitana I }\begin{center}
% latex table generated in R 3.3.1 by xtable 1.8-2 package
% Tue Sep  6 10:49:32 2016
\begin{tabular}{c|ccccccc}
  \hline
SE & temperatura & tweet & casos notif & casos preditos & ICmin & ICmax & incidência \\ 
  \hline
201630 & 17 & 0 & 181 & 200 & 194 & 201 & 2 \\ 
  201631 & 18 & 0 & 98 & 111 & 105 & 112 & 1 \\ 
  201632 & 18 & 0 & 93 & 113 & 107 & 114 & 1 \\ 
  201633 & 19 & 0 & 105 & 146 & 138 & 149 & 1 \\ 
  201634 & 16 & 0 & 70 & 121 & 109 & 124 & 1 \\ 
  201635 & 19 &  & 34 & 92 & 79 & 96 & 0 \\ 
   \hline
\end{tabular}

\end{center}
\small{\hyperlink{vartab}{ver descrição das variáveis}}\vspace{0.5cm}\begin{center}
\captionof{figure}{Casos notificados de dengue e Índice de menção em midia social sobre dengue na Regional Metropolitana I }\includegraphics[width=1.3\textwidth]{figs/tweetRJ_MetropolitanaI.png}\end{center}
\end{minipage}\hfill\begin{minipage}[t]{.30\linewidth}
\begin{mdframed}[style=sidebar,frametitle={}]
\textbf{\hyperlink{municips}{Municipios}}\begin{itemize}\gsquare Belford Roxo 
\gsquare Itaguaí 
\gsquare Magé 
\gsquare Mesquita 
\gsquare Queimados 
\gsquare Rio de Janeiro 
\gsquare São João de Meriti 
\gsquare Nova Iguaçu 
\gsquare Seropédica 
\gsquare Duque de Caxias 
\gsquare Japeri 
\gsquare Nilópolis 
\end{itemize}\BackToContents\end{mdframed}\hfill\end{minipage}\newpage\begin{minipage}[t]{.66\linewidth}
\hypertarget{Srrn}{\heading{Regional Serrana}{6pt}}
\includegraphics[width=0.8\textwidth]{figs/MapaRJ_Serrana.png}\vspace{0.5cm}\captionof{table}{Resumo das últimas seis semanas epidemiológicas na Regional Serrana }\begin{center}
% latex table generated in R 3.3.1 by xtable 1.8-2 package
% Tue Sep 27 12:40:36 2016
\begin{tabular}{c|ccccccc}
  \hline
SE & temperatura & tweet & casos notif & casos preditos & ICmin & ICmax & incidência \\ 
  \hline
201633 & 19 & 0 & 31 & 32 & 31 & 32 & 3 \\ 
  201634 & 16 & 0 & 57 & 61 & 58 & 62 & 6 \\ 
  201635 & 19 & 4 & 64 & 72 & 66 & 73 & 7 \\ 
  201636 & 20 & 0 & 5 & 5 & 5 & 5 & 1 \\ 
  201637 & 20 & 0 & 49 & 71 & 63 & 73 & 5 \\ 
  201638 & 19 & 0 & 0 & 0 & 0 & 0 & 0 \\ 
   \hline
\end{tabular}

\end{center}
\small{\hyperlink{vartab}{ver descrição das variáveis}}\vspace{0.5cm}\begin{center}
\captionof{figure}{Casos notificados de dengue e Índice de menção em midia social sobre dengue na Regional Serrana }\includegraphics[width=1.3\textwidth]{figs/tweetRJ_Serrana.png}\end{center}
\end{minipage}\hfill\begin{minipage}[t]{.30\linewidth}
\begin{mdframed}[style=sidebar,frametitle={}]
\textbf{\hyperlink{municips}{Municipios}}\begin{itemize}\gsquare Bom Jardim 
\gsquare Cachoeiras de Macacu 
\gsquare Cantagalo 
\gsquare Carmo 
\gsquare Duas Barras 
\gsquare Guapimirim 
\gsquare Nova Friburgo 
\gsquare Petrópolis 
\gsquare Santa Maria Madalena 
\gsquare São José do Vale do Rio Preto 
\gsquare Sumidouro 
\gsquare Teresópolis 
\gsquare Trajano de Moraes 
\gsquare São Sebastião do Alto 
\gsquare Cordeiro 
\gsquare Macuco 
\end{itemize}\BackToContents\end{mdframed}\hfill\end{minipage}\newpage\begin{minipage}[t]{.66\linewidth}
\hypertarget{Nort}{\heading{Regional Norte}{6pt}}
\includegraphics[width=0.8\textwidth]{figs/MapaRJ_Norte.png}\vspace{0.5cm}\captionof{table}{Resumo das últimas seis semanas epidemiológicas na Regional Norte }\begin{center}
% latex table generated in R 3.3.1 by xtable 1.8-2 package
% Tue Oct 25 11:50:13 2016
\begin{tabular}{c|ccccccc}
  \hline
SE & temperatura & tweet & casos notif & casos preditos & ICmin & ICmax & incidência \\ 
  \hline
201637 & 21 & 1 & 5 & 5 & 5 & 5 & 1 \\ 
  201638 & 20 & 1 & 10 & 10 & 10 & 10 & 1 \\ 
  201639 & 20 & 0 & 9 & 9 & 9 & 9 & 1 \\ 
  201640 & 21 & 2 & 11 & 13 & 11 & 13 & 1 \\ 
  201641 & 22 & 1 & 2 & 2 & 2 & 2 & 0 \\ 
  201642 & 24 &  & 0 & 0 & 0 & 0 & 0 \\ 
   \hline
\end{tabular}

\end{center}
\small{\hyperlink{vartab}{ver descrição das variáveis}}\vspace{0.5cm}\begin{center}
\captionof{figure}{Casos notificados de dengue e Índice de menção em midia social sobre dengue na Regional Norte }\includegraphics[width=1.3\textwidth]{figs/tweetRJ_Norte.png}\end{center}
\end{minipage}\hfill\begin{minipage}[t]{.30\linewidth}
\begin{mdframed}[style=sidebar,frametitle={}]
\textbf{\hyperlink{municips}{Municipios}}\begin{itemize}\gsquare Carapebus 
\gsquare Conceição de Macabu 
\gsquare Macaé 
\ysquare Quissamã 
\ysquare São Francisco de Itabapoana 
\ysquare São Fidélis 
\ysquare São João da Barra 
\ysquare Campos dos Goytacazes 
\end{itemize}\BackToContents\end{mdframed}\hfill\end{minipage}\newpage\begin{minipage}[t]{.66\linewidth}
\hypertarget{CntS}{\heading{Regional Centro Sul}{6pt}}
\includegraphics[width=0.8\textwidth]{figs/MapaRJ_CentroSul.png}\vspace{0.5cm}\captionof{table}{Resumo das últimas seis semanas epidemiológicas na Regional Centro Sul }\begin{center}
% latex table generated in R 3.3.1 by xtable 1.8-2 package
% Mon Sep 19 10:16:18 2016
\begin{tabular}{c|ccccccc}
  \hline
SE & temperatura & tweet & casos notif & casos preditos & ICmin & ICmax & incidência \\ 
  \hline
201631 & 18 & 0 & 0 & 0 & 0 & 0 & 0 \\ 
  201632 & 18 & 0 & 2 & 2 & 2 & 2 & 1 \\ 
  201633 & 19 & 0 & 0 & 0 & 0 & 0 & 0 \\ 
  201634 & 16 & 0 & 0 & 0 & 0 & 0 & 0 \\ 
  201635 & 19 &  & 0 & 0 & 0 & 0 & 0 \\ 
  201636 &  &  & 0 & 0 & 0 & 0 & 0 \\ 
   \hline
\end{tabular}

\end{center}
\small{\hyperlink{vartab}{ver descrição das variáveis}}\vspace{0.5cm}\begin{center}
\captionof{figure}{Casos notificados de dengue e Índice de menção em midia social sobre dengue na Regional Centro Sul }\includegraphics[width=1.3\textwidth]{figs/tweetRJ_CentroSul.png}\end{center}
\end{minipage}\hfill\begin{minipage}[t]{.30\linewidth}
\begin{mdframed}[style=sidebar,frametitle={}]
\textbf{\hyperlink{municips}{Municipios}}\begin{itemize}\gsquare Comendador Levy Gasparian 
\gsquare Engenheiro Paulo de Frontin 
\gsquare Mendes 
\gsquare Paracambi 
\gsquare Paraíba do Sul 
\gsquare Sapucaia 
\gsquare Três Rios 
\gsquare Vassouras 
\gsquare Paty do Alferes 
\gsquare Areal 
\gsquare Miguel Pereira 
\end{itemize}\BackToContents\end{mdframed}\hfill\end{minipage}\newpage\begin{minipage}[t]{.66\linewidth}
\hypertarget{MtII}{\heading{Regional Metropolitana II}{6pt}}
\includegraphics[width=0.8\textwidth]{figs/MapaRJ_MetropolitanaII.png}\vspace{0.5cm}\captionof{table}{Resumo das últimas seis semanas epidemiológicas na Regional Metropolitana II }\begin{center}
% latex table generated in R 3.3.1 by xtable 1.8-2 package
% Tue Sep 27 12:40:44 2016
\begin{tabular}{c|ccccccc}
  \hline
SE & temperatura & tweet & casos notif & casos preditos & ICmin & ICmax & incidência \\ 
  \hline
201633 & 20 & 2 & 28 & 30 & 28 & 30 & 1 \\ 
  201634 & 18 & 2 & 7 & 7 & 7 & 7 & 0 \\ 
  201635 & 21 & 3 & 11 & 11 & 11 & 13 & 1 \\ 
  201636 & 21 & 3 & 11 & 13 & 11 & 14 & 1 \\ 
  201637 & 22 & 1 & 11 & 15 & 12 & 17 & 1 \\ 
  201638 & 20 & 0 & 3 & 6 & 4 & 7 & 0 \\ 
   \hline
\end{tabular}

\end{center}
\small{\hyperlink{vartab}{ver descrição das variáveis}}\vspace{0.5cm}\begin{center}
\captionof{figure}{Casos notificados de dengue e Índice de menção em midia social sobre dengue na Regional Metropolitana II }\includegraphics[width=1.3\textwidth]{figs/tweetRJ_MetropolitanaII.png}\end{center}
\end{minipage}\hfill\begin{minipage}[t]{.30\linewidth}
\begin{mdframed}[style=sidebar,frametitle={}]
\textbf{\hyperlink{municips}{Municipios}}\begin{itemize}\gsquare Itaboraí 
\gsquare Niterói 
\gsquare Rio Bonito 
\gsquare Silva Jardim 
\gsquare São Gonçalo 
\gsquare Tanguá 
\gsquare Maricá 
\end{itemize}\BackToContents\end{mdframed}\hfill\end{minipage}\newpage

%-----------------------------------------------------------------------------------%	MAIN BODY - THIRD PAGE
%-----------------------------------------------------------------------------------
 %\begin{minipage}[t]{1\linewidth} % Mini page taking up 100% of the actual page

 
      \hypertarget{municips}{\heading{Resumo da situação epidemiológica da dengue nos municípios de Rio de Janeiro, na semana 43 de 2016}{6pt}}


     \begin{center}
            % latex table generated in R 3.3.1 by xtable 1.8-2 package
% Mon Sep 19 10:16:00 2016
\begin{longtable}{l|lllllll}
  \hline
Municipio & Regional & Temperatura & Tweets & Casos & Incidencia & Rt & Nivel \\ 
  \hline
\endhead
\hline
{\footnotesize Continua na próxima página}
\endfoot
\endlastfoot
Angra dos Reis & Litoral Sul &  &  & 9 & 5.4 & 0.8 & verde \\ 
  Aperibé & Noroeste & 22.1 &  & 0 & 0.0 & 0.0 & verde \\ 
  Araruama & Baixada Litorânea &  &  & 0 & 0.0 & 0.0 & verde \\ 
  Areal & Centro Sul & 20.6 &  & 0 & 0.0 & 0.0 & verde \\ 
  Armação dos Búzios & Baixada Litorânea &  &  & 0 & 0.0 & 0.0 & verde \\ 
  Arraial do Cabo & Baixada Litorânea & 21.7 &  & 0 & 0.0 & 0.0 & verde \\ 
  Barra do Piraí & Médio Paraíba &  &  & 0 & 0.0 & 0.0 & verde \\ 
  Barra Mansa & Médio Paraíba &  &  & 0 & 0.0 & 0.0 & verde \\ 
  Belford Roxo & Metropolitana I &  &  & 0 & 0.0 & 0.0 & verde \\ 
  Bom Jardim & Serrana &  &  & 0 & 0.0 & 0.0 & verde \\ 
  Bom Jesus do Itabapoana & Noroeste & 22.1 &  & 0 & 0.0 & 0.0 & verde \\ 
  Cabo Frio & Baixada Litorânea & 21.7 &  & 0 & 0.0 & 0.0 & verde \\ 
  Cachoeiras de Macacu & Serrana & 20.6 &  & 0 & 0.0 & 0.0 & verde \\ 
  Cambuci & Noroeste & 22.1 &  & 0 & 0.0 & 0.0 & verde \\ 
  Campos dos Goytacazes & Norte & 22.1 &  & 0 & 0.0 & 0.0 & verde \\ 
  Cantagalo & Serrana &  &  & 0 & 0.0 & 0.0 & verde \\ 
  Carapebus & Norte &  &  & 0 & 0.0 & 0.0 & verde \\ 
  Cardoso Moreira & Noroeste & 22.1 &  & 0 & 0.0 & 0.0 & verde \\ 
  Carmo & Serrana & 20.6 &  & 0 & 0.0 & 0.0 & verde \\ 
  Casimiro de Abreu & Baixada Litorânea & 21.7 &  & 0 & 0.0 & 0.0 & verde \\ 
  Comendador Levy Gasparian & Centro Sul & 20.6 &  & 0 & 0.0 & 0.0 & verde \\ 
  Conceição de Macabu & Norte &  &  & 1 & 4.5 & 3.1 & verde \\ 
  Cordeiro & Serrana &  &  & 0 & 0.0 & 0.0 & verde \\ 
  Duas Barras & Serrana &  &  & 0 & 0.0 & 0.0 & verde \\ 
  Duque de Caxias & Metropolitana I & 20.6 &  & 0 & 0.0 & 0.0 & verde \\ 
  Engenheiro Paulo de Frontin & Centro Sul & 20.6 &  & 0 & 0.0 & 0.0 & verde \\ 
  Guapimirim & Serrana & 21.3 &  & 0 & 0.0 & 0.0 & verde \\ 
  Iguaba Grande & Baixada Litorânea &  &  & 0 & 0.0 & 0.0 & verde \\ 
  Itaboraí & Metropolitana II & 21.3 &  & 0 & 0.0 & 0.0 & verde \\ 
  Itaguaí & Metropolitana I &  &  & 0 & 0.0 & 0.0 & verde \\ 
  Italva & Noroeste & 22.1 &  & 0 & 0.0 & 0.0 & verde \\ 
  Itaocara & Noroeste &  &  & 0 & 0.0 & 0.0 & verde \\ 
  Itaperuna & Noroeste & 22.1 &  & 0 & 0.0 & 0.0 & verde \\ 
  Itatiaia & Médio Paraíba &  &  & 0 & 0.0 & 0.0 & verde \\ 
  Japeri & Metropolitana I &  &  & 0 & 0.0 & 0.0 & verde \\ 
  Laje do Muriaé & Noroeste & 22.1 &  & 0 & 0.0 & 0.0 & verde \\ 
  Macaé & Norte &  &  & 5 & 4.8 & 1.0 & verde \\ 
  Macuco & Serrana &  &  & 0 & 0.0 & 0.0 & verde \\ 
  Magé & Metropolitana I & 20.6 &  & 0 & 0.0 & 0.0 & verde \\ 
  Mangaratiba & Litoral Sul &  &  & 0 & 0.0 & 0.0 & verde \\ 
  Maricá & Metropolitana II & 21.3 &  & 0 & 0.0 & 0.0 & verde \\ 
  Mendes & Centro Sul &  &  & 0 & 0.0 & 0.0 & verde \\ 
  Mesquita & Metropolitana I & 20.6 &  & 0 & 0.0 & 0.0 & verde \\ 
  Miguel Pereira & Centro Sul & 20.6 &  & 0 & 0.0 & 0.0 & verde \\ 
  Miracema & Noroeste & 22.1 &  & 0 & 0.0 & 0.0 & verde \\ 
  Natividade & Noroeste & 22.1 &  & 0 & 0.0 & 0.0 & verde \\ 
  Nilópolis & Metropolitana I &  &  & 0 & 0.0 & 0.0 & verde \\ 
  Niterói & Metropolitana II & 21.3 &  & 0 & 0.0 & 0.0 & verde \\ 
  Nova Friburgo & Serrana &  &  & 0 & 0.0 & 0.0 & verde \\ 
  Nova Iguaçu & Metropolitana I &  &  & 0 & 0.0 & 0.0 & verde \\ 
  Paracambi & Centro Sul &  &  & 0 & 0.0 & 0.0 & verde \\ 
  Paraíba do Sul & Centro Sul & 20.6 &  & 0 & 0.0 & 0.0 & verde \\ 
  Paraty & Litoral Sul &  &  & 1 & 2.5 & 0.5 & verde \\ 
  Paty do Alferes & Centro Sul & 20.6 &  & 0 & 0.0 & 0.0 & verde \\ 
  Petrópolis & Serrana & 20.6 &  & 0 & 0.0 & 0.0 & verde \\ 
  Pinheiral & Médio Paraíba &  &  & 0 & 0.0 & 0.0 & verde \\ 
  Piraí & Médio Paraíba &  &  & 0 & 0.0 & 0.0 & verde \\ 
  Porciúncula & Noroeste & 22.1 &  & 0 & 0.0 & 0.0 & verde \\ 
  Porto Real & Médio Paraíba &  &  & 0 & 0.0 & 0.0 & verde \\ 
  Quatis & Médio Paraíba &  &  & 0 & 0.0 & 0.0 & verde \\ 
  Queimados & Metropolitana I &  &  & 0 & 0.0 & 0.0 & verde \\ 
  Quissamã & Norte & 22.1 &  & 0 & 0.0 & 0.0 & verde \\ 
  Resende & Médio Paraíba &  &  & 0 & 0.0 & 0.0 & verde \\ 
  Rio Bonito & Metropolitana II &  &  & 0 & 0.0 & 0.0 & verde \\ 
  Rio Claro & Médio Paraíba &  &  & 0 & 0.0 & 0.0 & verde \\ 
  Rio das Flores & Médio Paraíba &  &  & 0 & 0.0 & 0.0 & verde \\ 
  Rio das Ostras & Baixada Litorânea &  &  & 0 & 0.0 & 0.0 & verde \\ 
  Rio de Janeiro & Metropolitana I &  &  & 25 & 1.1 & 0.7 & verde \\ 
  Santa Maria Madalena & Serrana &  &  & 0 & 0.0 & 0.0 & verde \\ 
  Santo Antônio de Pádua & Noroeste & 22.1 &  & 0 & 0.0 & 0.0 & verde \\ 
  São Fidélis & Norte & 22.1 &  & 0 & 0.0 & 0.0 & verde \\ 
  São Francisco de Itabapoana & Norte & 22.1 &  & 0 & 0.0 & 0.0 & verde \\ 
  São Gonçalo & Metropolitana II & 21.3 &  & 3 & 0.6 & 0.4 & verde \\ 
  São João da Barra & Norte & 22.1 &  & 0 & 0.0 & 0.0 & verde \\ 
  São João de Meriti & Metropolitana I & 21.3 &  & 0 & 0.0 & 0.0 & verde \\ 
  São José de Ubá & Noroeste & 22.1 &  & 0 & 0.0 & 0.0 & verde \\ 
  São José do Vale do Rio Preto & Serrana & 20.6 &  & 0 & 0.0 & 0.0 & verde \\ 
  São Pedro da Aldeia & Baixada Litorânea &  &  & 0 & 0.0 & 0.0 & verde \\ 
  São Sebastião do Alto & Serrana &  &  & 0 & 0.0 & 0.0 & verde \\ 
  Sapucaia & Centro Sul & 20.6 &  & 0 & 0.0 & 0.0 & verde \\ 
  Saquarema & Baixada Litorânea &  &  & 0 & 0.0 & 0.0 & verde \\ 
  Seropédica & Metropolitana I &  &  & 0 & 0.0 & 0.0 & verde \\ 
  Silva Jardim & Metropolitana II &  &  & 0 & 0.0 & 0.0 & verde \\ 
  Sumidouro & Serrana & 20.6 &  & 0 & 0.0 & 0.0 & verde \\ 
  Tanguá & Metropolitana II & 21.3 &  & 0 & 0.0 & 0.0 & verde \\ 
  Teresópolis & Serrana & 20.6 &  & 0 & 0.0 & 0.0 & verde \\ 
  Trajano de Moraes & Serrana &  &  & 0 & 0.0 & 0.0 & verde \\ 
  Três Rios & Centro Sul & 20.6 &  & 0 & 0.0 & 0.0 & verde \\ 
  Valença & Médio Paraíba &  &  & 0 & 0.0 & 0.0 & verde \\ 
  Varre-Sai & Noroeste & 22.1 &  & 0 & 0.0 & 0.0 & verde \\ 
  Vassouras & Centro Sul & 20.6 &  & 0 & 0.0 & 0.0 & verde \\ 
  Volta Redonda & Médio Paraíba &  &  & 0 & 0.0 & 0.0 & verde \\ 
  \hline
\end{longtable}

     \end{center}


      \BackToContents % Link back to the contents of the newsletter


\newpage

%---------------------------------------------------------------------------------
%	Variáves nas Tabelas, Créditos e Contato
%---------------------------------------------------------------------------------

\begin{minipage}[t]{1\linewidth} 

\hypertarget{vartab}{\heading{Lista das variáveis apresentadas nas tabelas:}{6pt}}

\begin{description}
\item [SE =] semana epidemiológica
\item [tweet =] número de tweets indicativos de casos de dengue na cidade
\item [temperatura =] média das temperaturas mínimas da semana
\item [casos notif =] casos notificados de dengue 
\item [casos preditos =] número de casos estimados após correção pelo atraso de notificação
\item [ICmin =] número mínimo de casos estimados (IC 95\%)
\item [ICmax =] número máximo de casos estimados (IC 95\%)
\item [Rt] número reprodutivo efetivo ($>$ 1 indica aumento de casos transmissão)
\item [p(Rt1) =] probabilidade do número reprodutivo ser maior que 1 ($>0.95$ indica aumento significativo de casos)
\item [inc =] incidência por 100.000 habitantes
\item [Nivel =] cor do alerta (verde, amarelo, laranja, vermelho)
\end{description}

\hypertarget{notas}{\heading{Notas}{6pt}}

\begin{itemize}
\item Os dados do sinan mais recentes ainda não foram totalmente digitados. Estimamos o número esperado de casos notificados considerando o tempo até os casos serem digitados.
\item Os dados de tweets são gerados pelo Observatório de Dengue (UFMG). Os tweets são processados para exclusão de informes e outros temas relacionados a dengue.
\item Algumas vezes, os casos da última semana ainda não estao disponíveis, nesse caso, usa-se uma estimação com base na tendência de variação da série.
\end{itemize}

\hypertarget{creditos}{\heading{Créditos}{6pt}}

Este é um projeto desenvolvido com apoio da SVS/MS em parceria com:

\begin{itemize}
\item Programa de Computação Científica, Fundação Oswaldo Cruz, Rio de Janeiro.
\item Escola de Matemática Aplicada, Fundação Getúlio Vargas.
\item Secretarias do Estado e Município do Rio de Janeiro.
\item Observatório de Dengue da UFMG
\item Secretaria Estadual de Saúde do Paraná.
\end{itemize}

      \BackToContents % Link back to the contents of the newsletter

\vspace{1cm}

\hrule
Para mais detalhes sobre o sistema de alerta InfoDengue, consultar: \url{http://info.dengue.mat.br}\\

\textbf{Contato}: \href{alerta\_dengue@fiocruz.br}{\nolinkurl{alerta\_dengue@fiocruz.br} }
\end{minipage} % fim da pagina de creditos

\end{document} 
